% Class-compatibility preamble retained for a reader excerpt of
%   Sahar Diskin, Anna Geisler, "Perfect Matching in Product Graphs and in their
%   Random Subgraphs", arXiv:2404.14020v4, CC BY 4.0.
%
% The source pins its own publisher class (scrartcl); the delivery wrapper accepts
% only the classes article and amsart, so this file supplies the amsart baseline and
% then copies the author's own declaration block mechanically and unchanged
% (pinned-original-source.tex, source lines 30-57: the author's newtheorem
% declarations and the author's macro declarations).  The author's own package
% list of source lines 4-21 is reused only for the entries the retained text needs
% (enumitem, amsmath, amssymb, amsthm, mathtools, xcolor, geometry) and the
% cross-reference packages already needed by the author's own \Cref/\eqref usage
% (hyperref, cleveref); tikz, caption, subcaption, float, graphicx, wrapfig and
% abstract serve figures and the title block that this excerpt does not carry.
% No mathematical macro, statement, symbol or assumption is introduced, changed
% or removed by this file.
\documentclass[10pt,a4paper]{amsart}
\usepackage[shortlabels]{enumitem}
\usepackage{amsmath}
\usepackage{amssymb}
\usepackage{amsthm}
\usepackage{mathtools}
\usepackage{xcolor}
\usepackage[margin=1in]{geometry}
\usepackage{hyperref}
\usepackage{cleveref}
\hypersetup{
  colorlinks=true,
  linkcolor=blue,
  filecolor=magenta,
  urlcolor=cyan
}

\newtheorem{theorem}{Theorem}
\newtheorem{thm}{Theorem}[section]
\newtheorem{corollary}[thm]{Corollary}
\newtheorem{lemma}[thm]{Lemma}
\newtheorem{proposition}[thm]{Proposition}
\newtheorem{claim}[thm]{Claim}
\newtheorem{remark}[thm]{Remark}
\newtheorem{question}[thm]{Question}

\newcommand{\whp}{whp}
\newcommand{\defi}{\operatorname{def}}
\newcommand{\odd}{\operatorname{odd}}
\newcommand{\iso}{\operatorname{iso}}
\newcommand{\cN}{\mathcal{N}}
\newcommand{\cO}{\mathcal{O}}
\newcommand{\cE}{\mathcal{E}}
\newcommand{\cK}{\mathcal{K}}
\newcommand{\cP}{\mathcal{P}}
\newcommand{\cQ}{\mathcal{Q}}

\newcommand{\footremember}[2]{%
  \footnote{#2}%
  \newcounter{#1}%
  \setcounter{#1}{\value{footnote}}%
}

\newcommand{\ag}[1]{\textcolor{purple}{(Anna: #1)}}
\newcommand{\sd}[1]{\textcolor{red}{(Sahar: #1)}}
\title{A connected d-regular graph with second-largest eigenvalue at least d-1 has a perfect or nearly-perfect matching}
\author{Sahar Diskin, Anna Geisler}
\date{arxiv:2404.14020v4; CC BY 4.0}
\begin{document}
\maketitle
\noindent\textit{Editorial note.} All delivered mathematical text is the author's own, copied from the pinned author source sources/combinatorics-fl04-vspace-1-5cm-perfect-matching-in-p/pinned-original-source.tex; no statement or proof step was rewritten, reordered or removed.\par\medskip
\noindent\textit{Editorial note.} Preamble provenance: the source pins its own publisher class scrartcl, which the delivery wrapper does not accept (supported wrapper classes are article and amsart), so this card uses sources/combinatorics-fl04-vspace-1-5cm-perfect-matching-in-p/retained-preamble.tex, a class-compatibility preamble that supplies the amsart baseline and then carries the author's own newtheorem and macro declaration block (pinned source lines 30-57) verbatim; its own header states which lines it copies and which of the author's packages were kept because the retained text needs them. No macro was invented and no mathematical symbol was changed. The host TeX installation was checked for every package that file loads and all of them are installed on the build host.\par\medskip
\noindent\textit{Editorial note.} Bibliography provenance: the counted statement and its proof cite nothing, so this card carries no bibliography at all. The source itself has no thebibliography environment and ends with bibliographystyle abbrv plus bibliography perc (author source lines 1691-1692); the author's attribution of the displayed context formula eq:Tutte-Berge to Bollobas is a sentence of the source (author source line 238) that the card does not count, and it is declared instead as a named standard prerequisite in the editorial list on the first page. No bibliography entry is reconstructed or invented.\par\medskip
\noindent\textit{Editorial note.} The displayed formula that the counted proof invokes, source label eq:Tutte-Berge, is delivered verbatim as a context span so that the closing step of the proof refers to a formula present in the assembled material; it is the author's own text and is not a replacement or a reconstruction.\par\medskip
\noindent\textit{Editorial note.} The author's source contains a commented-out proof environment for th:hitting-time at line 1646; it is a LaTeX comment in the author's file, carries no typeset content, and is not part of this card.\par\medskip
\noindent\textit{Allowed standard prerequisites (editorial).}\begin{itemize}
\item The Tutte-Berge formula for the number of unmatched vertices in a maximum matching. Used in the last line of the counted proof (source line 365, referring to the author's displayed formula eq:Tutte-Berge at source lines 240-241, itself attributed by the author to Bollobas at line 238). Converts the bound q <= |U| on the number of odd components of the graph minus U into the existence of a matching covering all vertices when the order is even and all but at most one vertex when the order is odd.
\end{itemize}\medskip
\begin{lemma}\label{lem:regular-matching}
Let $d \geq 1$ be an integer and let $\Gamma$ be a connected $d$-regular graph. Suppose that $\lambda(\Gamma)\ge d-1$.
Then $\Gamma$ has a perfect matching when $|V(\Gamma)|$ is even and a nearly-perfect matching when $|V(\Gamma)|$ is odd.
\end{lemma}

\begin{align} \label{eq:Tutte-Berge}
    \defi(\Gamma)= \max_{U\subseteq V(\Gamma)} \left\{\odd\left(\Gamma[V(\Gamma)\setminus U]\right)- |U|\right\},
\end{align}

\begin{proof}
For $U\subseteq V(\Gamma)$, let $Q_1,\ldots,Q_q$ be the odd components of $\Gamma-U$ such that $q=\odd(\Gamma - U)$.
If $U=\varnothing$, then $q=0$ when $|V(\Gamma)|$ is even and $q=1$ when it is odd.
Suppose now that $U\ne\varnothing$.
Every $Q_i$ is a proper non-empty set, so by the assumption that $\lambda(\Gamma)\geq d-1$,
\[
  e_\Gamma(Q_i,Q_i^c)\ge d-1.
\]
Since $|Q_i|$ is odd and $\Gamma$ is $d$-regular (implying that $d|Q_i|= 2 e_{\Gamma}(Q_i) + e_{\Gamma}(Q_i, Q_i^c)$), we have
\[
  e_\Gamma(Q_i,Q_i^c)\equiv d|Q_i|\equiv d\pmod 2.
\]
The cut has the same parity as $d$, whereas $d-1$ has the opposite parity; hence
\[
  e_\Gamma(Q_i,U)=e_\Gamma(Q_i,Q_i^c)\ge d.
\]
Summing over the odd components yields
\[
  qd\le\sum_{i=1}^q e_\Gamma(Q_i,U)\le d|U|,
\]
and therefore $q\le|U|$.
The Tutte--Berge formula \eqref{eq:Tutte-Berge} implies that all vertices are covered by a maximal matching if $|V(\Gamma)|$ is even and all but at most one vertex are covered by a maximal matching if $|V(\Gamma)|$ is odd.
\end{proof}

\end{document}
